% CHECKPOINT33: new source; preserves all checkpoint32 reconstruction bytes.
\statement{A.103}{Proved}{The uniform affine pentagon and the full signed congruence hexagon}
Let \(F=\widehat F(x,y)\), and let \(u\in F'\) satisfy duality
\(u(y,x)=u(x,y)^{-1}\) and the original marked sphere pentagon
\[
 u(a_0,c_0e_0)u(b_0c_0,f_0)
   =u(c_0,f_0)u(a_0b_0,e_0f_0)u(a_0,c_0).
 \tag{A.103.1}
\]
Here the six letters are the original
\((A_{12},A_{13},A_{23},A_{14},A_{24},A_{34})\) in the
completion of \(G=P_4/\langle\Delta_4^2\rangle\).
In the universal dihedral group
\(\widehat{\mathbb Z}_{T}\rtimes\langle s\rangle\), use
\(x\mapsto s,\ y\mapsto sT\), and define \(k\) by
\(u(s,sT)=T^k\). Put
\[
 d=1+2k,\qquad z=(xy)^{-1},\qquad
 H_d=u(x,y)x^k u(z,x)z^k u(y,z)y^k.
 \tag{A.103.2}
\]
Then \(k\in2\widehat{\mathbb Z}\), and \(d\) is a unit.
For the original signed representation
\[
 \begin{gathered}
 q:F\longrightarrow\operatorname{SL}_2(\widehat{\mathbb Z}),
 \qquad
 X=q(x)=\begin{pmatrix}1&2\\0&1\end{pmatrix},\quad
 Y=q(y)=\begin{pmatrix}1&0\\-2&1\end{pmatrix},\\
 U=q(u),\qquad b=U_{12},
 \end{gathered}
 \tag{A.103.3}
\]
one has the exact identities
\[
 dU_{11}=1,\qquad
 U=\begin{pmatrix}d^{-1}&b\\b&d(1+b^2)\end{pmatrix},
 \qquad q(H_d)=I.
 \tag{A.103.4}
\]
Every assertion holds in every signed congruence quotient, including
even and mixed moduli. For an actual canonical word of a negative
full-GT involution, A.37 supplies exactly (A.103.1).
The theorem itself requires no norm equation.

\begin{proof}
The closed derived subgroup of the universal dihedral group is
\(2\widehat{\mathbb Z}_{T}\): its commutators generate the even
rotations, and its quotient by those rotations is abelian.
This proves the asserted parity of \(k\).

\emph{The complete original source.}
Use the reversible A.96 presentation
\(G=F(P,A,C)\rtimes F(d_0,f_0)\), with positive conjugation actions
\[
 \begin{array}{c|ccc}
    &P&A&C\\ \hline
 d_0&C^{-1}PC&C^{-1}P^{-1}CPA&C\\
 f_0&A^{-1}PA&A&CPA^{-1}P^{-1}A
 \end{array}
 \tag{A.103.5}
\]
and inverse marking
\[
 a_0=Af_0,\quad c_0=Cd_0,\quad
 b_0=a_0^{-1}P^{-1}c_0^{-1},\quad
 e_0=d_0^{-1}Pf_0^{-1}.
 \tag{A.103.6}
\]
All six rows are equivalently the six relations
\[
 [A,f_0]=[C,d_0]=[Af_0,P]=[Cd_0,P]
       =[PA,d_0]=[CP,f_0]=1.
 \tag{A.103.7}
\]
Indeed the first two fix \(A,C\), the next two determine the
\(P\)-images, and the last two then give the remaining images
in (A.103.5). Thus no relation is omitted in using (A.103.7).
The boundary \(P\) is the original unrestricted generator.

\emph{A single uniform construction over every modulus.}
Fix a positive integer \(m\), put \(R=\mathbb Z/m\),
\(v=(1,0)^t,\ w=(0,1)^t\), and use
\((t,M)(s,N)=(t+Ms,MN)\) in
\(R^2\rtimes\operatorname{SL}_2(R)\).
The symbols \(X,Y,U,d,k\) in this paragraph denote their reductions.
Assign
\[
 \begin{gathered}
 P=(0,-I),\quad A=(-2v,I),\quad C=(-2w,I),\\
 f_0=(2v,X),\qquad d_0=(2w,Y).
 \end{gathered}
 \tag{A.103.8}
\]
These assignments satisfy the entire source. The first two
commutators in (A.103.7) vanish since \(Xv=v\) and \(Yw=w\).
The words \(Af_0=(0,X)\) and \(Cd_0=(0,Y)\) commute with \(P\).
For the last two commutators, the required translation identities are
\[
 (I-Y)(2v)=4w,\qquad (I-X)(-2w)=4v.
 \tag{A.103.9}
\]
They hold as polynomial identities over \(R\); nothing is divided
by two. Thus the homomorphism exists for even and nonreduced rings
as well, and extends continuously to \(\widehat G\).

The inverse marking gives
\[
 \begin{gathered}
 a_0=(0,X),\quad c_0=(0,Y),\\
 b_0=(0,-X^{-1}Y^{-1}),\quad
 e_0=(2v+2w,-Y^{-1}X^{-1}).
 \end{gathered}
 \tag{A.103.10}
\]
All five pairs in (A.103.1), in their original order, are
\[
 \begin{array}{c|c|c}
 \text{pair}&\text{first image}&\text{second image}\\ \hline
 (a_0,c_0e_0)&(0,X)&(2(v-w),-X^{-1})\\
 (b_0c_0,f_0)&(0,-X^{-1})&(2v,X)\\
 (c_0,f_0)&(0,Y)&(2v,X)\\
 (a_0b_0,e_0f_0)&(0,-Y^{-1})&(-2w,-Y^{-1})\\
 (a_0,c_0)&(0,X)&(0,Y).
 \end{array}
 \tag{A.103.11}
\]

\emph{The five evaluations.}
In each of the first two pair subgroups the linear heads commute
and preserve the direct summand \(Rv\). A commutator-word value
therefore has linear head \(I\). On the affine quotient by \(Rv\),
the first generator of pair one is the identity, and the second
generator of pair two is the identity. Each corresponding
\(u\)-value is consequently a translation in \(Rv\).
Only these pair subgroups are quotiented: \(Rv\) is not asserted
invariant under the entire affine target. The full left side of
(A.103.1) is thus \((\ell v,I)\) for some \(\ell\in R\).

For pair three put \(t=-w\).
Since \((I-Y)t=0\) and \((I-X)t=2v\), this is the simultaneous
translation conjugate of the linear pair \((Y,X)\).
Duality gives
\[
 u(c_0,f_0)=(t-U^{-1}t,U^{-1}).
 \tag{A.103.12}
\]
For pair four, the common decoration \(h=(0,Y^{-1})\) commutes
with both \(j_0=(0,-I)\) and \(j_1=(-2w,-I)\).
The pair is \((j_0h,j_1h)\). Central decorations disappear in a
commutator word because both exponent sums are zero; this holds
first for words and then for profinite words by continuity.
Both \(j_0,j_1\) are involutions and \(j_0j_1=(2w,I)\).
The actual universal dihedral definition of \(k\) gives
\[
 u(a_0b_0,e_0f_0)=(2kw,I).
 \tag{A.103.13}
\]
No exact-order assertion is made about this finite dihedral image,
so the argument includes all degenerate moduli.
The last value is \(u(a_0,c_0)=(0,U)\).
Multiplication in the displayed pentagon order now yields
\[
 (\ell v,I)=(-w+(1+2k)U^{-1}w,I).
 \tag{A.103.14}
\]
The second coordinate of \(U^{-1}w\) is \(U_{11}\), since
\(\det U=1\). Its vertical coordinate proves \(dU_{11}=1\) in
every \(R\). The compatible entries therefore satisfy that identity
in \(\widehat{\mathbb Z}\), proving that \(d\) is a full unit.

The integral matrix
\[
 S_0=\begin{pmatrix}0&-1\\1&0\end{pmatrix}
 \tag{A.103.15}
\]
interchanges \(X,Y\) by conjugation. Duality says
\(S_0US_0^{-1}=U^{-1}\). Comparing the two explicit matrix formulas
gives \(U_{12}=U_{21}\), with no division.
Together with the determinant and \(U_{11}=d^{-1}\), this proves
the complete matrix formula in (A.103.4).

\emph{The full signed hexagon, not only a selected modulus.}
Use the integral matrix
\[
 R_0=\begin{pmatrix}0&-1\\1&-1\end{pmatrix},\qquad
 R_0^3=I,\qquad Z=(XY)^{-1}.
 \tag{A.103.16}
\]
Direct multiplication over \(\mathbb Z\) gives the cycle
\(\operatorname{Ad}(R_0):X\mapsto Y\mapsto -Z\mapsto X\).
The original substitution
\(\rho:x\mapsto y,\ y\mapsto z\) differs, after \(q\), from
\(\operatorname{Ad}(R_0)\circ q\) by the central sign character
which is positive on \(x\) and negative on \(y\).
This character is trivial on \(F'\). Consequently, with
\(\kappa=\operatorname{Ad}(R_0)\),
\[
 q(u(y,z))=\kappa(U),\qquad
 q(u(z,x))=\kappa^2(U).
 \tag{A.103.17}
\]
The sign in \(Z^k\) disappears because \(k\) is even.
For \(V=UX^k\), the exact ordered product (A.103.2) is therefore
\[
 q(H_d)=V\kappa^2(V)\kappa(V)=(UX^kR_0^{-1})^3.
 \tag{A.103.18}
\]
This calculation uses the integral \(R_0\).
Its reduction modulo three equals \(XYX^{-1}\) there; no such
equality is asserted over \(\widehat{\mathbb Z}\).

Finally
\[
 X^kR_0^{-1}=\begin{pmatrix}-d&1\\-1&0\end{pmatrix},
 \qquad
 \operatorname{tr}(UX^kR_0^{-1})=-dU_{11}=-1.
 \tag{A.103.19}
\]
Its determinant is one. Cayley--Hamilton over the commutative ring
\(\widehat{\mathbb Z}\) gives \(g^2+g+I=0\), hence \(g^3=I\),
for \(g=UX^kR_0^{-1}\). Equation (A.103.18) proves the full
congruence assertion, including all prime powers and mixed moduli.

A.37 verifies the required full pentagon for the actual canonical
word in full GT. A.88 and A.100 already prove its unit and
regular-coefficient conclusions by the original Alexander route;
the present proof supplies the exact signed matrix and the full
signed congruence value without replacing those earlier proofs.
The assertion \(q(H_d)=I\) is not \(H_d=1\) in \(F\).
It does not prove \(\Omega_\alpha=1\), nor does it by itself
identify an independently normalized CM character.
A general GT0 canonical word is not granted a four-strand pentagon.
\end{proof}
